% Part: applied-modal-logics
% Chapter: epistemic-logic
% Section: introduction

\documentclass[../../../include/open-logic-section]{subfiles}

\begin{document}

\olfileid{aml}{el}{int}

\olsection{引言}

正如模态逻辑处理\emph{模态命题}及其间的后承关系，认知逻辑处理的是\emph{认知命题}及其间的后承关系。不再将模态算子解释为表示可能性与必然性，而是以认知或信念的方式解释一元联结词，从而刻画知识与信念。例如，我们可能想要表达下面这样的断言：

\begin{enumerate}
\item Richard 知道 Calgary 在 Alberta。
\item Audrey 认为狗可能在沙发上。
\item Richard 知道 Audrey 知道她的课在星期二。
\item 每个人都知道一年有 12 个月。
\end{enumerate}

当代认知逻辑常可追溯至 Jaako Hintikka（亚科·欣提卡）于 1962 年出版的《知识与信念》，该书成书于可能世界语义学在逻辑学中日益普及的时代。事实上，认知逻辑所使用的语义工具大多与其他模态逻辑相同，但对它们的解释有所不同。主要的变化在于我们将\emph{可达关系}用来表示什么。在认知逻辑中，可达关系代表某种形式的\emph{认知可能性}。我们将会看到，我们所刻画的认知概念会影响我们希望对可达关系施加的约束条件。我们也将看到，当对应理论被赋予认知解释时会发生什么。你会注意到，上面的例子提到了两个主体：Richard 和 Audrey，以及各自所知内容之间的关系。我们将要考察的认知逻辑是多主体逻辑，在其中可以表达此类关系。相反，单主体认知逻辑则只谈论单个个体的所知或所信。

\end{document}